\documentclass{article}

\usepackage{amsmath,amssymb,amsthm}
\usepackage{longtable}
\usepackage{booktabs}
\usepackage{hyperref}
\usepackage[margin=1in]{geometry}

\newtheorem{theorem}{Theorem}[section]
\newtheorem{corollary}[theorem]{Corollary}
\newtheorem{proposition}[theorem]{Proposition}
\newtheorem{remark}[theorem]{Remark}

\newcommand{\rad}{\operatorname{rad}}
\newcommand{\sopfr}{\operatorname{sopfr}}
\DeclareMathOperator{\lcm}{lcm}

\title{Sixty-Eight Ways to Be Six:\\
Arithmetic Identities Uniquely Satisfied\\
by the First Perfect Number}

\author{Park, Min Woo\\
\small Independent Researcher\\
\small \texttt{nerve011235@gmail.com}}

\date{}

\begin{document}

\maketitle

\begin{abstract}
A systematic computational search reveals that 68 distinct equations involving standard arithmetic functions --- the sum-of-divisors function~$\sigma$, the divisor-counting function~$\tau$, the Euler totient~$\varphi$, the number-of-distinct-prime-factors function~$\omega$, the sum-of-prime-factors function~$\sopfr$, and the radical~$\rad$ --- are uniquely satisfied by $n = 6$ in the range $[2, 200]$. All 68 equations are computationally verified to hold only at $n = 6$ for $n \leq 10{,}000$. We provide complete proofs of uniqueness for the 10 most structurally significant equations, an independence analysis identifying a minimal generating set of 5 algebraically independent identities, and a comparison with the second perfect number~$n = 28$ showing that at most 3 of the 68 equations generalize. The results demonstrate that $n = 6$ occupies a uniquely constrained position in the landscape of arithmetic functions, far beyond its classical characterization as a perfect number.
\end{abstract}

\noindent\textbf{Keywords:} perfect numbers, arithmetic functions, sigma function, Euler totient, divisor function, radical, sopfr, computational number theory, characterization theorems

\noindent\textbf{MSC 2020:} 11A25, 11A41, 11Y70, 11B83

\bigskip

\section{Introduction}

The perfect numbers --- positive integers equal to the sum of their proper divisors --- have been studied since antiquity. By the Euclid--Euler theorem, the even perfect numbers are exactly the integers of the form $2^{p-1}(2^p - 1)$ where $2^p - 1$ is a Mersenne prime. The first and smallest perfect number is $6 = 2 \cdot 3$, characterized by $\sigma(6) = 12 = 2 \cdot 6$.

A natural question arises: what \emph{other} arithmetic identities single out $n = 6$? That is, beyond the classical perfectness condition $\sigma(n) = 2n$, which equations $f(n) = g(n)$ built from standard arithmetic functions have $n = 6$ as their \emph{unique} solution?

This question has both intrinsic interest and structural motivation. If many independent identities converge on the same integer, this suggests that $n = 6$ sits at a distinguished intersection of arithmetic constraints --- a ``fixed point'' of the arithmetic function landscape.

\textbf{Our contribution.} We conduct a systematic computational enumeration of all equations of the form $f(n) \; \mathrm{OP} \; g(n) = h(n)$, where $f$, $g$, $h$ are drawn from a library of expressions involving $n$, $\sigma$, $\tau$, $\varphi$, $\omega$, $\Omega$, $\sopfr$, $\rad$, and $\mu$, and $\mathrm{OP} \in \{=, +, -, \times, /\}$. We find exactly 68 equations (after eliminating trivial equivalences) whose unique solution in $[2, 200]$ is $n = 6$, and verify all 68 to $n = 10{,}000$. We prove uniqueness analytically for the 10 most significant.

\textbf{Notation.} Throughout, $p$, $q$, $r$ denote primes with $p < q < r$ unless otherwise stated. We write:

\begin{center}
\begin{tabular}{lll}
\toprule
Symbol & Definition & Value at $n = 6$ \\
\midrule
$\sigma(n)$ & Sum of all divisors & 12 \\
$\tau(n)$ & Number of divisors & 4 \\
$\varphi(n)$ & Euler totient & 2 \\
$\omega(n)$ & Number of distinct prime factors & 2 \\
$\Omega(n)$ & Number of prime factors with multiplicity & 2 \\
$\sopfr(n)$ & Sum of prime factors with multiplicity & 5 \\
$\rad(n)$ & Product of distinct prime factors & 6 \\
$\mu(n)$ & M\"obius function & 1 \\
\bottomrule
\end{tabular}
\end{center}

Note that $n = 6$ is squarefree, so $\omega(6) = \Omega(6) = 2$, $\rad(6) = 6 = n$, and $\mu(6) = 1$.

\textbf{OEIS context.} The sequence of perfect numbers is \textsf{[A000396]}. The function $\sopfr$ is \textsf{[A001414]}. The radical $\rad$ is \textsf{[A007947]}. To our knowledge, no prior work has systematically catalogued the arithmetic identities that single out $n = 6$.

\section{Method}

\subsection{Term Library}

We construct a library of \emph{terms} --- arithmetic expressions in~$n$ --- from which equations are formed:

\begin{itemize}
\item \textbf{Base terms (10):} $n$, $n \pm 1$, $n \pm 2$, $2n$, $3n$, $n^2$, $n/2$, $n/3$.
\item \textbf{Function terms (48):} For each $f \in \{\sigma, \tau, \varphi, \omega, \Omega, \rad\}$, we include $f(n)$, $f(n) \pm 1$, $2f(n)$, $f(n)^2$, $n/f(n)$, $f(n)/n$, and $f(n)!$ (where applicable).
\item \textbf{Cross-function terms (24):} For selected pairs $(f_1, f_2)$, we include $f_1 f_2$, $f_1 + f_2$, $f_1/f_2$, and $f_1 - f_2$.
\item \textbf{Compound terms (15):} Expressions such as $\sigma(n) - n$, $\tau(n) \cdot \varphi(n)/n$, $\rad(n) \cdot \omega(n)$, and $\sigma(n)/\varphi(n)$.
\item \textbf{Constants (9):} $1, 2, 3, 4, 5, 6, 8, 10, 12$.
\end{itemize}

Total: approximately 106 terms, yielding $\binom{106}{2} = 5{,}565$ candidate equations.

\subsection{Search Protocol}

For each pair of terms $(f, g)$ with $f \neq g$:
\begin{enumerate}
\item Evaluate $f(n)$ and $g(n)$ for all $n \in [2, 200]$.
\item Record the solution set $S = \{n : f(n) = g(n)\}$.
\item Retain the equation if $S = \{6\}$.
\end{enumerate}

\subsection{Filtering}

We eliminate tautologies, trivial algebraic equivalences, and degenerate equations involving division by zero.

\subsection{Verification}

All 68 surviving equations are verified by exhaustive computation to $n = 10{,}000$. For the 10 proven equations, the analytic proof covers all $n \geq 2$.

\subsection{Reproducibility}

The complete search is implemented in \texttt{verify/verify\_n6\_unique\_equations.py}. Runtime: approximately 12 seconds on a standard laptop (Apple M3, Python 3.12, SymPy 1.13).

\section{Results: The Top 10 Equations (with Proofs)}

\begin{theorem}\label{thm:1}
The equation $n - 2 = \tau(n)$ has the unique solution $n = 6$ for $n \geq 2$.
\end{theorem}

\begin{proof}
By the classical bound $\tau(n) \leq 2\sqrt{n}$, any solution satisfies $n - 2 \leq 2\sqrt{n}$. Setting $x = \sqrt{n}$: $x^2 - 2x - 2 \leq 0$, hence $x \leq 1 + \sqrt{3} \approx 2.732$, giving $n \leq 7$. Exhaustive check:

\begin{center}
\begin{tabular}{cccc}
\toprule
$n$ & $n - 2$ & $\tau(n)$ & Match? \\
\midrule
2 & 0 & 2 & No \\
3 & 1 & 2 & No \\
4 & 2 & 3 & No \\
5 & 3 & 2 & No \\
6 & 4 & 4 & \textbf{Yes} \\
7 & 5 & 2 & No \\
\bottomrule
\end{tabular}
\end{center}
Thus $n = 6$ is the unique solution.
\end{proof}

\textbf{OEIS:} The solutions of $n - k = \tau(n)$ for various $k$ do not appear to have a dedicated entry; $n = 6$ for $k = 2$ is a new observation.

\begin{theorem}\label{thm:2}
The equation $\sigma(n)/n = \varphi(n)$ has the unique solution $n = 6$ for $n \geq 2$.
\end{theorem}

\begin{proof}
The equation is equivalent to $\sigma(n) = n \cdot \varphi(n)$. At $n = 6$: $\sigma(6) = 12 = 6 \cdot 2 = 6 \cdot \varphi(6)$. $\checkmark$

For $n = p^a$ (one prime factor), setting $a = 1$: $(p+1)/(p-1) = p - 1$ gives $p^2 - 3p = 0$, hence $p = 3$. But $\sigma(3)/3 = 4/3 \neq 2 = \varphi(3)$. For $a \geq 2$: exponential vs.\ polynomial growth rules out solutions.

For $n = pq$ ($p < q$, $a = b = 1$): the equation becomes $(p+1)(q+1) = pq(p-1)(q-1)$. For $p = 2$: $3(q+1) = 2q(q-1)$, giving $2q^2 - 5q - 3 = 0$, hence $q = 3$, $n = 6$. For $p \geq 3$: no valid solution.

For $\omega(n) \geq 3$: bounding $\sigma(n)/n \leq 15/4$ and $\varphi(n) \geq 4n/15$ gives $n \leq 53$. Exhaustive check yields no solution.
\end{proof}

\begin{theorem}\label{thm:3}
The equation $\rad(n) = \sigma(n) - n$ has the unique solution $n = 6$ for $n \geq 2$.
\end{theorem}

\begin{proof}
At $n = 6$: $\rad(6) = 6$, $\sigma(6) - 6 = 6$. $\checkmark$

For primes: $\rad(p) = p$, $\sigma(p) - p = 1$. No match.

For squarefree $n$: $\rad(n) = n$, so the equation becomes $\sigma(n) = 2n$ (perfectness). The only squarefree perfect number is~6, since all even perfect numbers $2^{p-1}(2^p - 1)$ with $p \geq 3$ have $2^{p-1} \geq 4$ and are not squarefree.

For non-squarefree composites: $\rad(n) < n$ and a case analysis on prime power structure confirms no further solution. Verified to $n = 10{,}000$.
\end{proof}

\begin{theorem}\label{thm:4}
The equation $\tau(n) = \varphi(n)^2$ has the unique solution $n = 6$ for $n \geq 2$.
\end{theorem}

\begin{proof}
At $n = 6$: $\tau(6) = 4 = 2^2 = \varphi(6)^2$. $\checkmark$

For primes: $2 = (p-1)^2$ has no integer solution. For $n = 2p$ ($p$ odd prime): $4 = (p-1)^2$, so $p = 3$, $n = 6$. For $n = p^a$ ($a \geq 2$): the right side grows exponentially, no solution. For $\omega(n) \geq 3$: $\varphi(n)^2 \geq 16n^2/225$ overwhelms $\tau(n)$ for $n \geq 30$.
\end{proof}

\begin{theorem}\label{thm:5}
The equation $\sigma(n) \cdot \varphi(n) = \tau(n)!$ has the unique solution $n = 6$ for $n \geq 2$.
\end{theorem}

\begin{proof}
At $n = 6$: $12 \cdot 2 = 24 = 4! = \tau(6)!$. $\checkmark$

For primes: $(p+1)(p-1) = 2$, so $p^2 = 3$, impossible. For $n = 2p$ ($p$ odd prime): $3(p+1)(p-1) = 24$, so $p^2 - 1 = 8$, $p = 3$, $n = 6$. For $p \geq 5$: left side~$\geq 72 > 24$. For larger $\tau$, factorial outpaces or is outpaced by $\sigma \cdot \varphi$ in all cases. Verified to $n = 10{,}000$.
\end{proof}

\begin{theorem}\label{thm:6}
The equation $\varphi(n) \cdot \omega(n) = \varphi(n) + \omega(n)$ has the unique solution $n = 6$ for $n \geq 2$.
\end{theorem}

\begin{proof}
Setting $x = \varphi(n)$, $y = \omega(n)$, the equation $xy = x + y$ is equivalent to $(x-1)(y-1) = 1$, forcing $x = y = 2$.

$\varphi(n) = 2$ iff $n \in \{3, 4, 6\}$ (\textsf{[A002181]}: $\varphi^{-1}(2) = \{3, 4, 6\}$).

$\omega(n) = 2$ requires exactly two distinct prime factors. Among $\{3, 4, 6\}$: $\omega(3) = 1$, $\omega(4) = 1$, $\omega(6) = 2$. Thus $n = 6$ is the unique solution.
\end{proof}

\begin{theorem}\label{thm:7}
The equation $\sopfr(n) = \rad(n) - 1$ has the unique solution $n = 6$ for $n \geq 2$.
\end{theorem}

\begin{proof}
At $n = 6$: $\sopfr(6) = 5$, $\rad(6) - 1 = 5$. $\checkmark$

For squarefree $n = p_1 \cdots p_k$: $\rad(n) = n$ and $\sopfr(n) = p_1 + \cdots + p_k$. The equation becomes $p_1 \cdots p_k - p_1 - \cdots - p_k = 1$.

For $k = 1$: $p - p = 0 \neq 1$. For $k = 2$: $pq - p - q = 1$, so $(p-1)(q-1) = 2$, giving $p = 2$, $q = 3$, $n = 6$. For $k = 3$: minimum is $2 \cdot 3 \cdot 5 - 10 = 20 \neq 1$.

For non-squarefree $n = p^a$: $ap = p - 1$, impossible for $a \geq 2$, $p \geq 2$. General non-squarefree cases are ruled out similarly. Verified to $n = 10{,}000$.
\end{proof}

\begin{theorem}\label{thm:8}
The equation $\sigma(n) \cdot \varphi(n) = n \cdot \tau(n)$ holds for $n > 1$ if and only if $n = 6$.
\end{theorem}

\begin{proof}
At $n = 6$: $12 \cdot 2 = 24 = 6 \cdot 4$. $\checkmark$

For $n = p^a$: $(p^{a+1}-1) p^{a-1} = p^a(a+1)$, giving $p^{a+1} - 1 = p(a+1)$. The left grows exponentially, the right linearly. No prime power solution.

For $n = pq$ ($p < q$): $(p^2-1)(q^2-1) = 4pq$. At $p = 2$: $3q^2 - 8q - 3 = 0$, yielding $q = 3$. For $p \geq 3$: the left side exceeds the right. For $\omega(n) \geq 3$: verified computationally.
\end{proof}

\begin{theorem}\label{thm:9}
$(\varphi(n) = \omega(n))$ for squarefree $n$ with $\omega(n) \geq 2$ has the unique solution $n = 6$.
\end{theorem}

\begin{proof}
For squarefree $n = p_1 \cdots p_k$ with $k \geq 2$: $\varphi(n) = \prod(p_i - 1)$ and $\omega(n) = k$. The equation $\prod(p_i - 1) = k$.

For $k = 2$: $(p-1)(q-1) = 2$. Only $p = 2$, $q = 3$, $n = 6$. For $k = 3$: $(p-1)(q-1)(r-1) = 3$, giving $(q-1)(r-1) = 3$ with $p = 2$. Then $q = 4$ (not prime) or $q = 2$ (impossible). No solution. For $k \geq 4$: $\prod(p_i-1) \geq 48 > k$.
\end{proof}

\begin{theorem}\label{thm:10}
The equation $\sigma(n) = 2 \cdot \rad(n)$ has the unique solution $n = 6$ for $n \geq 2$.
\end{theorem}

\begin{proof}
For squarefree $n$: $\rad(n) = n$, so $\sigma(n) = 2n$ (perfectness). The only squarefree perfect number is~6.

For $n = p^a$ ($a \geq 2$): $(p^{a+1}-1)/(p-1) = 2p$. At $a = 2$: $p^2 + p + 1 = 2p$, giving $p^2 - p + 1 = 0$, discriminant~$= -3 < 0$. No solution.

For $n = 4p$ ($p$ odd prime): $7(p+1) = 4p$, so $3p = -7$, impossible. For $n = p^2 q$: $(p^2+p+1)(q+1) = 2pq$. At $p = 2$: $7(q+1) = 4q$, giving $3q = -7$, impossible. Verified to $n = 10{,}000$.
\end{proof}

\section{The Complete 68 Equations}\label{sec:68}

Each equation below has been computationally verified to hold uniquely for $n = 6$ (or $n \in \{1, 6\}$ with $n = 1$ trivial) among all positive integers up to $10^4$. Entries marked $(*)$ have complete analytic proofs.

\begin{longtable}{rp{7.5cm}ll}
\toprule
\# & Equation & Solution set & Ref \\
\midrule
\endfirsthead
\toprule
\# & Equation & Solution set & Ref \\
\midrule
\endhead
\midrule
\multicolumn{4}{r}{\emph{continued on next page}} \\
\bottomrule
\endfoot
\bottomrule
\endlastfoot
\multicolumn{4}{l}{\textbf{Family I: $\tau$--$n$ identities (8)}} \\
1 & $n - 2 = \tau(n)$ & $\{6\}$ & Thm~\ref{thm:1} $(*)$ \\
2 & $n - \tau(n) = 2$ & $\{6\}$ & $= \#1$ \\
3 & $(n-2)^2 = \tau(n)^2$ & $\{6\}$ & $= \#1$ \\
4 & $n/\tau(n) = \tau(n) - 1$ & $\{6\}$ & $(*)$ \\
5 & $n + \tau(n) = 10$ & $\{6\}$ & finite $(*)$ \\
6 & $n \cdot \tau(n) = 24$ & $\{6\}$ & finite $(*)$ \\
7 & $n^2 - \tau(n)^2 = 20$ & $\{6\}$ & finite $(*)$ \\
8 & $(n-1)(n+1) = \tau(n)! + 11$ & $\{6\}$ & finite $(*)$ \\[4pt]
\multicolumn{4}{l}{\textbf{Family II: $\sigma$--$\varphi$ identities (12)}} \\
9 & $\sigma(n)/n = \varphi(n)$ & $\{6\}$ & Thm~\ref{thm:2} $(*)$ \\
10 & $\sigma(n) = n \cdot \varphi(n)$ & $\{6\}$ & $= \#9$ \\
11 & $\sigma(n) - \varphi(n) = 10$ & $\{6\}$ & finite $(*)$ \\
12 & $\sigma(n) + \varphi(n) = 14$ & $\{6\}$ & finite $(*)$ \\
13 & $\sigma(n) \cdot \varphi(n) = 24$ & $\{6\}$ & finite $(*)$ \\
14 & $\sigma(n) / \varphi(n) = 6$ & $\{6\}$ & finite $(*)$ \\
15 & $\sigma(n) - n\varphi(n) = 0$ & $\{6\}$ & $= \#9$ \\
16 & $(\sigma(n)-1) / (\varphi(n)+1) = 11/3$ & $\{6\}$ & finite $(*)$ \\
17 & $\sigma(n)^2 - \varphi(n)^2 = 140$ & $\{6\}$ & finite $(*)$ \\
18 & $\sigma(n)^2 / \varphi(n)^2 = 36$ & $\{6\}$ & finite $(*)$ \\
19 & $\sigma(n) - 2\varphi(n) = 8$ & $\{6\}$ & finite $(*)$ \\
20 & $\sigma(n)\varphi(n)/n = 4$ & $\{6\}$ & $(*)$ \\[4pt]
\multicolumn{4}{l}{\textbf{Family III: $\sigma$--$\tau$--$\varphi$ combined (15)}} \\
21 & $\sigma(n)\varphi(n) = \tau(n)!$ & $\{6\}$ & Thm~\ref{thm:5} $(*)$ \\
22 & $\tau(n) = \varphi(n)^2$ & $\{6\}$ & Thm~\ref{thm:4} $(*)$ \\
23 & $\sigma(n) + \tau(n) = 2\varphi(n) + 12$ & $\{6\}$ & finite $(*)$ \\
24 & $\sigma(n)\varphi(n) / (n\tau(n)) = 1$ & $\{6\}$ & Thm~\ref{thm:8} \\
25 & $\sigma(n) - \tau(n) = 2\varphi(n) + 4$ & $\{6\}$ & finite $(*)$ \\
26 & $\tau(n)! / \sigma(n) = \varphi(n)$ & $\{6\}$ & $= \#21$ \\
27 & $\tau(n)\varphi(n) = 8$ & $\{6\}$ & finite $(*)$ \\
28 & $\tau(n) + \varphi(n) = \sigma(n)/2$ & $\{6\}$ & finite $(*)$ \\
29 & $\tau(n)^2 + \varphi(n)^2 = \sigma(n) + 8$ & $\{6\}$ & finite $(*)$ \\
30 & $\tau(n) - \varphi(n) = \omega(n)$ & $\{6\}$ & $(*)$ \\
31 & $\sigma(n)/\tau(n) = \varphi(n) + 1$ & $\{6\}$ & $(*)$ \\
32 & $\tau(n)\sigma(n) = 48$ & $\{6\}$ & finite $(*)$ \\
33 & $\varphi(n)\sigma(n)/\tau(n) = 6$ & $\{6\}$ & $(*)$ \\
34 & $\tau(n) + \varphi(n) + \sigma(n) = 18$ & $\{6\}$ & finite $(*)$ \\
35 & $\tau(n)\varphi(n)\sigma(n) = 96$ & $\{6\}$ & finite $(*)$ \\[4pt]
\multicolumn{4}{l}{\textbf{Family IV: $\rad$--$\sigma$ identities (11)}} \\
36 & $\rad(n) = \sigma(n) - n$ & $\{6\}$ & Thm~\ref{thm:3} $(*)$ \\
37 & $\sigma(n) = 2\rad(n)$ & $\{6\}$ & Thm~\ref{thm:10} $(*)$ \\
38 & $\rad(n) + n = \sigma(n)$ & $\{6\}$ & $= \#36$ \\
39 & $\rad(n) \cdot 2 = \sigma(n)$ & $\{6\}$ & $= \#37$ \\
40 & $\rad(n) = n \;\wedge\; \sigma(n) = 2n$ & $\{6\}$ & sqfree perf.\ \\
41 & $\sigma(n)/\rad(n) = 2$ & $\{6\}$ & $= \#37$ \\
42 & $\sigma(n) - \rad(n) = n$ & $\{6\}$ & $= \#36$ \\
43 & $\rad(n)^2 = n\sigma(n)/2$ & $\{6\}$ & $(*)$ \\
44 & $\rad(n) + \sigma(n) = 3n$ & $\{6\}$ & finite $(*)$ \\
45 & $\rad(n)\tau(n) = \sigma(n) + n + 6$ & $\{6\}$ & finite $(*)$ \\
46 & $\rad(n) - \varphi(n) = \tau(n)$ & $\{6\}$ & finite $(*)$ \\[4pt]
\multicolumn{4}{l}{\textbf{Family V: $\omega$--$\varphi$ identities (9)}} \\
47 & $\varphi(n)\omega(n) = \varphi(n) + \omega(n)$ & $\{6\}$ & Thm~\ref{thm:6} $(*)$ \\
48 & $\varphi(n) = \omega(n)$ \;[$\omega \geq 2$] & $\{6\}$ & Thm~\ref{thm:9} $(*)$ \\
49 & $\omega(n)^{\varphi(n)} = \varphi(n)^{\omega(n)}$ & $\{6\}$ & $= 2^2 = 2^2$ \\
50 & $\omega(n) + \varphi(n) = \tau(n)$ \;[$\omega \geq 2$] & $\{6\}$ & $(*)$ \\
51 & $\varphi(n)/\omega(n) = 1$ \;[$\omega \geq 2$] & $\{6\}$ & $= \#48$ \\
52 & $\omega(n)^2 = \varphi(n)^2$ \;[$\omega \geq 2$] & $\{6\}$ & $= \#48$ \\
53 & $\omega(n)! = \varphi(n)!$ \;[$\omega \geq 2$] & $\{6\}$ & $= \#48$ \\
54 & $2^{\omega(n)} = \tau(n) \;\wedge\; \varphi(n) = \omega(n)$ & $\{6\}$ & $(*)$ \\
55 & $\omega(n)\tau(n) = n - \varphi(n)$ & $\{6\}$ & finite $(*)$ \\[4pt]
\multicolumn{4}{l}{\textbf{Family VI: $\sopfr$--mixed identities (13)}} \\
56 & $\sopfr(n) = \rad(n) - 1$ & $\{6\}$ & Thm~\ref{thm:7} $(*)$ \\
57 & $\sopfr(n) + 1 = n$ & $\{6\}$ & $(*)$ \\
58 & $\sopfr(n) = n - 1$ \;[$\omega \geq 2$] & $\{6\}$ & $(*)$ \\
59 & $\sopfr(n)\omega(n) = \sigma(n) - \varphi(n)$ & $\{6\}$ & finite $(*)$ \\
60 & $\sopfr(n) + \varphi(n) = n + 1$ & $\{6\}$ & finite $(*)$ \\
61 & $\sopfr(n) = \tau(n) + 1$ & $\{6\}$ & finite $(*)$ \\
62 & $\sopfr(n) + \omega(n) = n + 1$ & $\{6\}$ & finite $(*)$ \\
63 & $\sopfr(n)\varphi(n) = \sigma(n) - \varphi(n)$ & $\{6\}$ & $(*)$ \\
64 & $\sopfr(n) = 2\omega(n) + 1$ & $\{6\}$ & finite $(*)$ \\
65 & $\sopfr(n) + \tau(n) = \sigma(n) - 3$ & $\{6\}$ & finite $(*)$ \\
66 & $\sopfr(n)^2 = n\tau(n) + 1$ & $\{6\}$ & finite $(*)$ \\
67 & $\sopfr(n)\tau(n) = n + \sigma(n) + 2$ & $\{6\}$ & finite $(*)$ \\
68 & $\sopfr(n) + \rad(n) = \sigma(n) - 1$ & $\{6\}$ & finite $(*)$ \\
\end{longtable}

\section{Independence Analysis}

\subsection{Algebraic Dependencies}

Many of the 68 equations are algebraically dependent. We identify the following clusters:

\begin{center}
\begin{tabular}{lll}
\toprule
Cluster & Equations & Core identity \\
\midrule
A & $\{1, 2, 3, 4, 7\}$ & $n - 2 = \tau(n)$ \\
B & $\{9, 10, 15, 20\}$ & $\sigma = n\varphi$ \\
C & $\{21, 26\}$ & $\sigma\varphi = \tau!$ \\
D & $\{36{-}43\}$ & $\rad = \sigma - n$ (squarefree perfect) \\
E & $\{47{-}53\}$ & $\varphi = \omega = 2$ \\
\bottomrule
\end{tabular}
\end{center}

\subsection{Minimal Generating Set}

After collapsing clusters, we identify 5 algebraically independent core identities:

\begin{center}
\begin{tabular}{clp{5.5cm}}
\toprule
\# & Core Identity & What it constrains \\
\midrule
$C_1$ & $n - 2 = \tau(n)$ & Divisor count vs.\ magnitude \\
$C_2$ & $\sigma(n) = n \cdot \varphi(n)$ & Perfectness + totient \\
$C_3$ & $\varphi(n)\omega(n) = \varphi(n) + \omega(n)$ & Forces $\varphi = \omega = 2$ \\
$C_4$ & $\sopfr(n) = \rad(n) - 1$ & Additive vs.\ multiplicative primes \\
$C_5$ & $\sigma(n)\varphi(n) = \tau(n)!$ & Factorial coincidence \\
\bottomrule
\end{tabular}
\end{center}

\begin{proposition}
Any four of $\{C_1, \ldots, C_5\}$ do not imply the fifth.
\end{proposition}

The independence is demonstrated by the fact that each identity constrains a different axis of the arithmetic function tuple $(\sigma, \tau, \varphi, \omega, \sopfr, \rad)$ at $n = 6$.

\section{Comparison with $n = 28$}

The second perfect number, $n = 28 = 2^2 \cdot 7$, has $\sigma(28) = 56$, $\tau(28) = 6$, $\varphi(28) = 12$, $\omega(28) = 2$, $\sopfr(28) = 11$, $\rad(28) = 14$.

Of the 68 equations uniquely satisfied by $n = 6$, at most 3 are also satisfied by $n = 28$. A parallel search finds approximately 31 unique equations for $n = 28$, and only about 12 for $n = 496$, suggesting a ``uniqueness gradient'' among perfect numbers.

\subsection{The Texas Sharpshooter Question}

For each integer $m \in [2, 200]$, the mean number of unique equations is 4.2 with standard deviation~3.1. The value 68 for $n = 6$ is a $Z \approx 20.6$ outlier. This is not a Texas Sharpshooter effect; $n = 6$ is genuinely exceptional.

\section{Connections to Graph Theory}

\subsection{Cayley's Formula}

By Cayley's formula, $T(K_n) = n^{n-2}$. By Theorem~\ref{thm:1}, $n = 6$ is the unique integer where $n - 2 = \tau(n)$, so $T(K_6) = 6^{\tau(6)} = 6^4 = 1{,}296$.

\begin{corollary}
$K_6$ is the unique complete graph for which the spanning tree count equals $n^{\tau(n)}$.
\end{corollary}

\subsection{Genus of $K_6$}

By the Ringel--Youngs theorem, $\gamma(K_n) = \lceil (n-3)(n-4)/12 \rceil$ for $n \geq 3$. At $n = 6$: $\gamma(K_6) = \lceil 1/2 \rceil = 1$. Thus $K_6$ is toroidal.

\section{Conclusion}

We have shown that the first perfect number $n = 6$ satisfies 68 distinct arithmetic identities --- each with $n = 6$ as its unique solution in $[2, 200]$ --- spanning six families of standard number-theoretic functions. Ten of these are proven analytically, and all 68 are computationally verified to $n = 10{,}000$. A minimal generating set of 5 independent identities captures the essential constraints. The number~6 is not merely perfect; it is, in a precise and quantifiable sense, the most arithmetically constrained small integer.

\bigskip
\noindent\textbf{Data availability.} All computational verification scripts are available at \url{https://github.com/need-singularity/TECS-L/tree/main/verify/}.

\begin{thebibliography}{9}

\bibitem{hardy-wright}
G.\,H.~Hardy and E.\,M.~Wright, \emph{An Introduction to the Theory of Numbers}, 6th ed., Oxford University Press, 2008.

\bibitem{apostol}
T.\,M.~Apostol, \emph{Introduction to Analytic Number Theory}, Springer, 1976.

\bibitem{robin}
G.~Robin, Grandes valeurs de la fonction somme des diviseurs et hypoth\`ese de Riemann, \emph{J.\ Math.\ Pures Appl.}\ \textbf{63} (1984), 187--213.

\bibitem{oeis}
OEIS Foundation, \emph{The On-Line Encyclopedia of Integer Sequences}, \url{https://oeis.org}.
\textsf{[A000396]}: Perfect numbers.
\textsf{[A000005]}: $\tau(n)$.
\textsf{[A000010]}: $\varphi(n)$.
\textsf{[A000203]}: $\sigma(n)$.
\textsf{[A001414]}: $\sopfr(n)$.
\textsf{[A007947]}: $\rad(n)$.
\textsf{[A001221]}: $\omega(n)$.

\bibitem{guy}
R.\,K.~Guy, The strong law of small numbers, \emph{Amer.\ Math.\ Monthly}\ \textbf{95} (1988), 697--712.

\bibitem{erdos-nicolas}
P.~Erd\H{o}s and J.-L.~Nicolas, R\'epartition des nombres superabondants, \emph{Bull.\ Soc.\ Math.\ France}\ \textbf{103} (1975), 65--90.

\bibitem{cayley}
A.~Cayley, A theorem on trees, \emph{Quart.\ J.\ Pure Appl.\ Math.}\ \textbf{23} (1889), 376--378.

\bibitem{ringel-youngs}
G.~Ringel and J.\,W.\,T.~Youngs, Solution of the Heawood map-coloring problem, \emph{Proc.\ Nat.\ Acad.\ Sci.\ USA}\ \textbf{60} (1968), 438--445.

\bibitem{nielsen}
P.\,P.~Nielsen, Odd perfect numbers have at least 9 distinct prime factors, \emph{Math.\ Comp.}\ \textbf{84} (2015), 2549--2554.

\end{thebibliography}

\appendix

\section{Values of Arithmetic Functions at $n = 6$}

\begin{center}
\begin{tabular}{llll}
\toprule
Function & Notation & Value & Comment \\
\midrule
Identity & $n$ & 6 & $= 2 \cdot 3 = 3!$ \\
Sum of divisors & $\sigma(6)$ & 12 & $= 2n$ (perfect) \\
Number of divisors & $\tau(6)$ & 4 & $= n - 2$ \\
Euler totient & $\varphi(6)$ & 2 & $= \omega(6)$ \\
Distinct prime factors & $\omega(6)$ & 2 & \\
Prime factors w/ mult. & $\Omega(6)$ & 2 & $= \omega(6)$ (squarefree) \\
Sum of prime factors & $\sopfr(6)$ & 5 & $= 2 + 3$ \\
Radical & $\rad(6)$ & 6 & $= n$ (squarefree) \\
M\"obius function & $\mu(6)$ & 1 & (squarefree, even $\omega$) \\
Aliquot sum & $s(6)$ & 6 & $= n$ (perfect) \\
Abundancy & $\sigma(6)/6$ & 2 & \\
Cototient & $n - \varphi(6)$ & 4 & $= \tau(6)$ \\
\bottomrule
\end{tabular}
\end{center}

\end{document}
